% Options for packages loaded elsewhere
\PassOptionsToPackage{unicode}{hyperref}
\PassOptionsToPackage{hyphens}{url}
\PassOptionsToPackage{dvipsnames,svgnames,x11names}{xcolor}
%
\documentclass[
  11pt,
]{article}
\usepackage{amsmath,amssymb}
\usepackage{iftex}
\ifPDFTeX
  \usepackage[T1]{fontenc}
  \usepackage[utf8]{inputenc}
  \usepackage{textcomp} % provide euro and other symbols
\else % if luatex or xetex
  \usepackage{unicode-math} % this also loads fontspec
  \defaultfontfeatures{Scale=MatchLowercase}
  \defaultfontfeatures[\rmfamily]{Ligatures=TeX,Scale=1}
\fi
\usepackage{lmodern}
\ifPDFTeX\else
  % xetex/luatex font selection
\fi
% Use upquote if available, for straight quotes in verbatim environments
\IfFileExists{upquote.sty}{\usepackage{upquote}}{}
\IfFileExists{microtype.sty}{% use microtype if available
  \usepackage[]{microtype}
  \UseMicrotypeSet[protrusion]{basicmath} % disable protrusion for tt fonts
}{}
\makeatletter
\@ifundefined{KOMAClassName}{% if non-KOMA class
  \IfFileExists{parskip.sty}{%
    \usepackage{parskip}
  }{% else
    \setlength{\parindent}{0pt}
    \setlength{\parskip}{6pt plus 2pt minus 1pt}}
}{% if KOMA class
  \KOMAoptions{parskip=half}}
\makeatother
\usepackage{xcolor}
\usepackage[margin=1in]{geometry}
\setlength{\emergencystretch}{3em} % prevent overfull lines
\providecommand{\tightlist}{%
  \setlength{\itemsep}{0pt}\setlength{\parskip}{0pt}}
\setcounter{secnumdepth}{-\maxdimen} % remove section numbering
\DeclareUnicodeCharacter{220E}{\ensuremath{\square}}
\ifLuaTeX
  \usepackage{selnolig}  % disable illegal ligatures
\fi
\IfFileExists{bookmark.sty}{\usepackage{bookmark}}{\usepackage{hyperref}}
\IfFileExists{xurl.sty}{\usepackage{xurl}}{} % add URL line breaks if available
\urlstyle{same}
\hypersetup{
  pdftitle={Typed Evidence Licenses for Finite Positive Rule Graphs},
  pdfauthor={Sze Chun Yiu},
  pdfsubject={Typed least-fixed-point semantics for finite positive scientific rule graphs},
  pdfkeywords={automated reasoning, scientific evidence, provenance, belief revision},
  colorlinks=true,
  linkcolor={Maroon},
  filecolor={Maroon},
  citecolor={Blue},
  urlcolor={Blue},
  pdfcreator={LaTeX via pandoc}}

\title{Typed Evidence Licenses for Finite Positive Rule Graphs}
\author{Sze Chun Yiu\\\texttt{sze-chun.yiu@fysik.su.se}}
\date{25 August 2026}

\begin{document}
\maketitle

\begin{abstract}
Boolean dependency graphs can report whether a claim remains reachable
after refutation, but not which evidence type licenses the surviving
derivation. We define a finite typed semantics for positive conjunctive
scientific rule graphs. Independent seeds carry declared licenses; every
rule has a license cap; and directly refuted claims receive the empty
label. The induced monotone operator has a least fixed point on a finite
powerset lattice.

We prove finite convergence, rule-order independence, and a proof-tree
characterization: a license reaches a claim exactly when a finite
untainted proof tree carries it through every seed and rule cap.
Unsupported cycles stay empty, added refutations can only remove
licenses, and the retraction contains exactly the claim-license pairs
that lose every untainted proof tree. Caps also prevent promotion:
post-outcome repairs cannot regain prospective authority, and bounded
computation cannot acquire a \(\textsf{THEOREM}\) license merely because
an untyped conclusion remains reachable.

A deterministic evaluator implements the semantics, while a JSON schema
validates document shape. Three bounded cases illustrate forecast
falsification, query-specific information falsification, and
nonpromotion of a computational frontier. The component is restricted to
finite positive rule graphs; it does not model negation, probability,
inconsistency, or general scientific judgment.

\end{abstract}

\noindent\textbf{Keywords:} automated reasoning; scientific evidence; provenance;
least fixed points; belief revision; executable semantics

\hypertarget{introduction}{%
\section{1. Introduction}\label{introduction}}

Scientific records mix claims supported by different evidence types:
analytic proof, constructive bounds, finite exact computation,
prospective prediction, forecast-only evidence, post-outcome repair, and
bounded computation awaiting external replay. A counterexample can
invalidate one layer while leaving another intact.

An untyped dependency graph can preserve independent derivations, but it
can still overpromote a surviving repair. If a post-outcome predictor
becomes exact on the observed panel, it should not inherit prospective
status merely because the old predictor and the repair share a
conclusion string. Similarly, a bounded exact search should not acquire
a theorem-grade license by passing through a rule whose conclusion is a
theorem-shaped sentence.

We attach evidence licenses to the least fixed point of a positive rule
graph. The substrate stays close to positive Datalog and provenance:
finite claims, positive conjunctive rules, monotone iteration, and
finite proof trees. The scientific addition is explicit nonpromotion.
Derivability and evidence license must travel together.

Our contributions are:

\begin{enumerate}
\def\labelenumi{\arabic{enumi}.}
\tightlist
\item
  powerset labels for typed evidence licenses on finite claims;
\item
  capped conjunctive transfer, which admits a license only when every
  premise and the rule cap admit it;
\item
  a proof-tree characterization of the least fixed point;
\item
  monotone and canonical semantic retraction under direct refutation;
\item
  a deterministic reusable evaluator and machine-readable schema; and
\item
  three bounded cases that expose the difference between reachability
  and licensed scientific use.
\end{enumerate}

\hypertarget{claims-licenses-and-rules}{%
\section{2. Claims, licenses, and
rules}\label{claims-licenses-and-rules}}

Let \(Q\) be a finite claim set and \(\Lambda\) a finite license set.
Example licenses are

\[
\begin{gathered}
\textsf{THEOREM}, \quad \textsf{CONSTRUCTIVE\_BOUND}, \quad
\textsf{FINITE\_EXACT}, \\
\textsf{PROSPECTIVE}, \quad \textsf{FORECAST\_ONLY}, \quad
\textsf{POST\_OUTCOME}, \\
\textsf{BOUNDED\_COMPUTATION}, \quad \textsf{EXTERNAL\_REPLAY}.
\end{gathered}
\]

A label is a subset of \(\Lambda\), ordered by inclusion. Its join is
union and its bottom element is the empty set. Each claim \(q\) has an
independent seed label \(\sigma(q)\subseteq\Lambda\).

A positive rule is a triple

\[
r=(A\to h, K_r),
\]

where \(A\subseteq Q\) is a nonempty finite body, \(h\in Q\) is the
head, and \(K_r\subseteq\Lambda\) is the rule cap. For premise labels
\((\ell_a)_{a\in A}\), define

\[
\tau_r((\ell_a)_{a\in A})
=K_r\cap\bigcap_{a\in A}\ell_a.
\]

A license crosses a rule only when every premise carries it and the cap
permits it. Rules with empty bodies are represented as independent
seeds. Let \(R\subseteq Q\) be the directly refuted claims.

\hypertarget{least-fixed-point-semantics}{%
\section{3. Least-fixed-point
semantics}\label{least-fixed-point-semantics}}

For a label assignment \(x\in(2^\Lambda)^Q\), define

\[
F_R(x)_q=
\begin{cases}
\varnothing,&q\in R, \\
\displaystyle\sigma(q)\cup
\bigcup_{\substack{r\,\mid\,\operatorname{head}(r)=q}}
\tau_r(x|_{\operatorname{body}(r)}),&q\notin R.
\end{cases}
\]

The operator is monotone. Starting from all-empty labels, iterate to
stabilization and denote the result by

\[
\operatorname{Lic}(R)=\operatorname{lfp}(F_R).
\]

Call an asynchronous evaluation \emph{accumulating} if it initializes
every unrefuted claim \(q\) with \(\sigma(q)\), initializes every
refuted claim with the empty label, and thereafter only unions a
nonempty rule transfer into an unrefuted rule head. A rule instance is
\emph{enabled} for a license \(\lambda\) when every body claim currently
carries \(\lambda\), the rule cap contains \(\lambda\), the head is
unrefuted, and the head does not yet carry \(\lambda\). It is
\emph{fair} if every enabled instance is eventually fired; after each
label growth, all rule instances are reconsidered.

\textbf{Theorem 1 (finite convergence and order independence).}
Synchronous bottom-up iteration has at most \(|Q||\Lambda|\) changing
rounds, followed by one stability check. Every fair accumulating rule
schedule reaches the same least fixed point.

\textbf{Proof.} With \(R\) fixed, labels can only gain licenses. There
are at most \(|Q||\Lambda|\) claim-license pairs. Under a fair
accumulating schedule, each pair having a finite proof tree is
eventually added by induction on tree height; no unsupported pair can be
added because every firing is an instance of the defining operator. The
stable assignment is therefore exactly the least fixed point,
independently of schedule. ∎

\hypertarget{typed-proof-trees}{%
\section{4. Typed proof trees}\label{typed-proof-trees}}

A proof tree for \((q, \lambda)\) is valid under \(R\) when:

\begin{enumerate}
\def\labelenumi{\arabic{enumi}.}
\tightlist
\item
  its root is \(q\), and no node is directly refuted;
\item
  a leaf \(a\) satisfies \(\lambda\in\sigma(a)\); and
\item
  an internal node applies a declared rule \(A\to h\) whose cap contains
  \(\lambda\), with one child proof tree carrying \(\lambda\) for every
  antecedent in \(A\).
\end{enumerate}

\textbf{Theorem 2 (proof-tree equivalence).} A license \(\lambda\)
belongs to \(\operatorname{Lic}(R)_q\) if and only if a finite valid
proof tree exists for \((q, \lambda)\).

\textbf{Proof.} Induction on the first iteration round in which
\(\lambda\) enters the label of \(q\) constructs a tree. Conversely,
induction on tree height shows that each leaf license enters from a seed
and crosses every internal rule through its cap. ∎

Discarding labels and keeping only claims with nonempty labels recovers
the ordinary reachability view, but loses the evidence distinction.

\hypertarget{cycles-and-nonpromotion}{%
\section{5. Cycles and nonpromotion}\label{cycles-and-nonpromotion}}

Rules \(a\to b\) and \(b\to a\) with no seed labels have the all-empty
least fixed point. A cycle supplies a derivation shape, not evidence. If
\(a\) has a theorem seed and both caps permit \(\textsf{THEOREM}\), that
license propagates to \(b\). If either cap excludes it, the license
stops at that edge.

\textbf{Corollary 3 (license conservation).} Every license at a
conclusion occurs in every leaf seed and every rule cap along at least
one finite proof tree.

A repair rule whose cap includes \(\textsf{POST\_OUTCOME}\) but excludes
\(\textsf{PROSPECTIVE}\) cannot transmit prospective authority, even
when one of its premises happens to carry both. Nonpromotion is
therefore enforced by the rule semantics rather than by an informal note
attached after evaluation.

\hypertarget{retraction-under-refutation}{%
\section{6. Retraction under
refutation}\label{retraction-under-refutation}}

\textbf{Theorem 4 (refutation monotonicity).} If \(R\subseteq R'\), then

\[
\operatorname{Lic}(R')_q\subseteq\operatorname{Lic}(R)_q
\]

for every claim \(q\).

\textbf{Proof.} For every \(x\), \(F_{R'}(x)\) is pointwise contained in
\(F_R(x)\). Iteration from bottom preserves containment. ∎

Let \(L_{\mathrm{pre}}=\operatorname{Lic}(\varnothing)\) and
\(L_{\mathrm{post}}=\operatorname{Lic}(R)\). Define

\[
\operatorname{Ret}(R)
=\{(q, \lambda):
\lambda\in L_{\mathrm{pre}}(q)\setminus L_{\mathrm{post}}(q)\}.
\]

Call a post-refutation assignment \emph{proof-supported} when it retains
every claim-license pair with a finite untainted proof tree and retains
no pair without such a tree. Order retractions by set inclusion.

\textbf{Theorem 5 (canonical semantic retraction).} Relative to the
declared seeds, rules, caps, and refutations, \(L_{\mathrm{post}}\) is
the unique proof-supported assignment. Consequently,
\(\operatorname{Ret}(R)\) removes every and only pair that loses all
finite untainted proof trees.

\textbf{Proof.} By Theorem 2, membership in \(L_{\mathrm{post}}\) is
equivalent to the existence of a finite valid proof tree under \(R\).
Any proof-supported assignment must therefore equal
\(L_{\mathrm{post}}\), and its complement relative to
\(L_{\mathrm{pre}}\) must equal \(\operatorname{Ret}(R)\). ∎

This is a semantic retraction relative to the declared system. It makes
no claim that the input license policy is the only reasonable scientific
policy.

\hypertarget{executable-semantics}{%
\section{7. Executable semantics}\label{executable-semantics}}

The public JSON Schema first validates document shape and required
fields. The reference evaluator then performs semantic validation: claim
and rule identifiers must be unique, and every license, rule-body claim,
rule head, and refutation must refer to a declared object. After both
layers pass, the evaluator performs deterministic bottom-up iteration.
Output records the final label of every claim, the number of iterations,
and the typed retraction from the unrefuted baseline.

The implementation rejects undeclared claims or licenses, empty rule
bodies, duplicate rule identifiers, and malformed caps. Canonical
sorting makes equal inputs produce byte-stable semantic outputs. The
evaluator is deliberately small enough for line-by-line comparison with
\(F_R\); it does not infer caps or scientific policy.

\hypertarget{bounded-case-encodings}{%
\section{8. Bounded case encodings}\label{bounded-case-encodings}}

The following are synthetic fixtures used to test the typed semantics.
Each fixture is self-contained and supplies its own seeds, rules, caps,
and refutations.

\hypertarget{forecast-falsification}{%
\subsection{8.1 Forecast falsification}\label{forecast-falsification}}

A synthetic optimization record contains an explicit feasible
construction, an independent all-size support theorem, and a compact
equality forecast supported by a forecast-labeled seed set. A held-out
synthetic result satisfies \(C_{\mathrm{exact}}=10<11\), directly
refuting the equality and its regime label. The construction and the
independent support theorem retain their licenses.

A repaired support-two forecast can carry a theorem-grade license
inherited from the independent theorem and a post-outcome license from
its construction. Its repair cap excludes \(\textsf{PROSPECTIVE}\), so
it cannot be counted as prospective confirmation. The old and repaired
panels remain distinct evidence objects.

\hypertarget{decision-survives-value-and-witness-falsifiers}{%
\subsection{8.2 Decision survives value and witness
falsifiers}\label{decision-survives-value-and-witness-falsifiers}}

A second self-contained fixture seeds a four-index certificate for unary
optimality. Its declared counterexamples show that complete pair
information does not determine exact improvement or the presence of a
triple block, and that all proper interaction marginals do not determine
exact value.

Decision, value, and optimizer-structure claims occupy separate nodes.
Refuting pair-value sufficiency removes its licenses without touching
the independently seeded decision theorem. This is query typing: one
representation can carry a theorem license for a decision query and no
exact-value license.

\hypertarget{a-bounded-support-frontier-does-not-decide-an-exact-constant}{%
\subsection{8.3 A bounded support frontier does not decide an exact
constant}\label{a-bounded-support-frontier-does-not-decide-an-exact-constant}}

A third self-contained fixture assigns analytic licenses to a width-one
generalized-Davenport corridor and a saturation-defect lemma. An exact
bounded search excludes a specified obstruction through a declared
finite support frontier, but its evidence contract labels the result as
bounded computation awaiting external replay.

Any rule from that frontier to a support claim is capped by the same
licenses. No proof tree carries \(\textsf{THEOREM}\) to either exact
candidate value of the generalized Davenport constant or to the
associated extremal classification. Repeated implementations by the same
research group do not become independent mathematical replication.

\hypertarget{relation-to-prior-work}{%
\section{9. Relation to prior work}\label{relation-to-prior-work}}

Truth-maintenance systems and belief revision own dependency-directed
update {[}1,2{]}. Positive Datalog owns least-fixed-point rule semantics
{[}3,4{]}. Annotated and semiring provenance owns typed query
annotations, alternative derivations, recursion, trust annotations, and
deletion behavior {[}5,6{]}. Work on recursive-query causality relates
minimal supports, causes, responsibility, and deletion robustness
{[}7,8{]}.

We claim no generic novelty for fixed points, proof trees, provenance
labels, minimal supports, hitting sets, causality, or deletion
robustness. The residual is a narrowly scoped scientific component: a
finite evidence-license vocabulary, cap-preserving nonpromotion, exact
typed retraction, a deterministic evaluator, and bounded cases in which
the license distinction changes what may be reported.

\hypertarget{reproducibility-and-limitations}{%
\section{10. Reproducibility and
limitations}\label{reproducibility-and-limitations}}

The evaluator is checked against exhaustive fixed-point enumeration on
bounded random rule graphs, and unit tests cover unsupported and seeded
cycles, cap blocking, alternative derivations, and refutation
monotonicity. The proofs carry the general finite authority.

The component handles positive conjunctive rules only. Stratified
negation, defaults, probabilistic evidence, and inconsistency are out of
scope. The license set and caps are curated policy inputs, not inferred
truths. Powerset intersection is one transparent transfer choice. The
cases demonstrate machine-checkable behavior in specified records; they
are not evidence of broad human-science usability.

\hypertarget{conclusion}{%
\section{11. Conclusion}\label{conclusion}}

Reachability alone is not scientific authority. A claim can remain
reachable under a weaker license, and a repaired claim can regain
exactness without regaining prospective status. The typed least fixed
point makes those distinctions explicit and executable.

The component fails closed in two directions: unsupported cycles create
no licenses, and underlicensed derivations create no stronger evidence
types. Its retraction is canonical relative to the declared inputs: it
removes exactly the pairs that lose every valid proof tree. No separate
inclusion-minimality claim is made.

\hypertarget{tool-use-disclosure}{%
\section{Tool-use disclosure}\label{tool-use-disclosure}}

A generative language model assisted manuscript organization, language
revision, and submission-package preparation. The listed author remains
responsible for the mathematical statements, proofs, references,
executable claims, and final text.

\hypertarget{data-and-code-availability}{%
\section{Data and code availability}\label{data-and-code-availability}}

The submission package includes the JSON schema, deterministic Python
evaluator, unit tests, and bounded case fixtures required to reproduce
every executable claim. These files are distributed in the source
archive accompanying this version.

\hypertarget{references}{%
\section{References}\label{references}}

\begin{enumerate}
\def\labelenumi{\arabic{enumi}.}
\tightlist
\item
  J. Doyle, ``A Truth Maintenance System,'' \emph{Artificial
  Intelligence} \textbf{12}, 231-272 (1979). DOI:
  10.1016/0004-3702(79)90008-0
\item
  J. P. Martins and S. C. Shapiro, ``A Model for Belief Revision,''
  \emph{Artificial Intelligence} \textbf{35}, 25-79 (1988). DOI:
  10.1016/0004-3702(88)90031-8
\item
  C. Bourgaux, P. Bourhis, L. Peterfreund, and M. Thomazo, ``Revisiting
  Semiring Provenance for Datalog,'' in \emph{KR 2022} (2022). DOI:
  10.24963/kr.2022/10
\item
  M. Abo Khamis, H. Q. Ngo, R. Pichler, D. Suciu, and Y. R. Wang,
  ``Convergence of Datalog over (Pre-)Semirings,'' in \emph{Proceedings
  of the 41st ACM SIGMOD-SIGACT-SIGAI Symposium on Principles of
  Database Systems} (PODS 2022), 105-117 (2022). DOI:
  10.1145/3517804.3524140
\item
  T. J. Green, G. Karvounarakis, and V. Tannen, ``Provenance
  Semirings,'' in \emph{Proceedings of PODS 2007}, 31-40 (2007). DOI:
  10.1145/1265530.1265535
\item
  P. A. Bonatti, A. Hogan, A. Polleres, and L. Sauro, ``Robust and
  Scalable Linked Data Reasoning Incorporating Provenance and Trust
  Annotations,'' \emph{Journal of Web Semantics} \textbf{9}(2), 165-201
  (2011). DOI: 10.1016/j.websem.2011.06.003
\item
  R. B. Thapa and S. Staab, ``Causality and Minimal Supports in
  Recursive Datalog,'' arXiv:2607.16443 (2026).
\item
  R. B. Thapa and S. Staab, ``Causal Explanations for Stratified
  Datalog,'' arXiv:2608.21141 (2026).
\end{enumerate}

\end{document}
